\section{与同期工作的关系}\label{app:comparison}

Qiushi Engine 独立形成了本文的结构性证明。Wang 的同期预印本建立了相同的数值下界\cite{wang2026twentyone}。我们在研究及科学报告完成之后的发布准备阶段，才获悉该工作。

下面对照 Wang 已公开的二元域证明与复现包，以及本文公开的证明与研究记录，区分共同的数学基础与有限前提、推导和研究贡献上的差异。

\paragraph{共同的数学基础。}
两种证明都继承了经过认证的限制下界，并使用占据容量；都导出了高秩第一因子两两之差秩一，并将它们置于共同的仿射行陪集或列陪集中。这些是实质性的共同归约，而不仅是术语相似。按零化子限制输入，与对系数空间取商，是同一运算在对偶约定下的表述。D'Ambrosio 的工作更早使用了容量论证与紧限制\cite{dambrosio2026}。

\begin{longtable}{@{}L{2.1cm}L{5.7cm}L{6.4cm}@{}}
\toprule
方面&Wang 的二元域证明&Qiushi Engine 的结构性证明\\
\midrule
\endfirsthead
\toprule 方面&Wang 的二元域证明&Qiushi Engine 的结构性证明\\\midrule
\endhead
\bottomrule\endfoot
有限前提
&通过秩一张成搜索与回溯加强限制表。
&以固定的早期表为基础，新增八个商下界，由整数占据编码与 CNF/DRAT 反驳证明认证。\\
陪集归约之后
&分类 35 种高秩构型，并依据加强后的表检验其秩一补全。
&将仿射平面容量与完整张量的秩和结合，强迫得到唯一秩型 $(16,1,3)$，并在秩和 27 处取等。\\
决定性推导
&通过带容量约束的穷尽枚举排除所有第一因子秩型，无须补全第二、第三因子。
&取等给出 $M_tP^{-1}M_s=\delta_{ts}M_t$；张量乘积公式将三个因子耦合起来，与必需的可逆性矛盾。\\
验证
&检查限制证书与完整秩型枚举。
&复验继承证书，重建八个占据系统，验证其 DRAT 证明，并在报告中给出后续代数推导。\\
进一步结果
&另证明了 $R_{\mathbb F_3}(\langle2,3,3\rangle)=15$。
&饱和引理还约束零超额的秩 22 分解：至多一个第一因子可逆，该位置只剩 $(17,5,0)$ 与 $(18,3,1)$ 两种秩型。\\
\end{longtable}

因此，区别不只是替换证明的最后一行：产生有限前提的计算不同，对共同陪集归约的使用方式不同，最终矛盾也从第一因子容量不可行性转为跨因子代数矛盾。共同的数值结论并不使这些机制相同。

\paragraph{研究发展与开放材料。}
Wang 提供了包括代码和证书的可复现定理包。我们的发布还记录了结构性证明如何形成：秩 22 构造与形变、商核心、支撑与补全检验、纠错、回到完整张量，以及对饱和条件的识别。附录\ref{app:research-record}将这一发展整合为与证据相联的 Meta-Trace，并由主题研究档案支持。在结构性论证之外，这也构成对自主数学研究进行记录和研究的贡献。

对秩一张成搜索、对称性约化的秩型枚举或三元域矩形情形感兴趣的读者，可以在 Wang 的文章\cite{wang2026twentyone}中找到互补方法。结合本文的饱和论证与 Meta-Trace 阅读，该工作为理解有限域约束如何导出张量秩下界，提供了另一种视角。
